\documentclass[10pt,a4paper]{amsart}
\usepackage{amsmath,amssymb,amsthm,graphicx}
\usepackage{verbatim}
\newenvironment{reference}{\comment}{\endcomment}
%\newenvironment{reference}{}{}
\newenvironment{slogan}{\comment}{\endcomment}
\newenvironment{history}{\comment}{\endcomment}

% For commutative diagrams we use Xy-pic
\usepackage[all]{xy}

% We use 2cell for 2-commutative diagrams.
\xyoption{2cell}
% Theorem environments.
%
\theoremstyle{plain}
\newtheorem{theorem}[subsection]{Theorem}
\newtheorem{proposition}[subsection]{Proposition}
\newtheorem{lemma}[subsection]{Lemma}

\theoremstyle{definition}
\newtheorem{definition}[subsection]{Definition}
\newtheorem{example}[subsection]{Example}
\newtheorem{exercise}[subsection]{Exercise}
\newtheorem{situation}[subsection]{Situation}

\theoremstyle{remark}
\newtheorem{remark}[subsection]{Remark}
\newtheorem{remarks}[subsection]{Remarks}

\numberwithin{equation}{subsection}

% Macros
%
\def\lim{\mathop{\mathrm{lim}}\nolimits}
\def\colim{\mathop{\mathrm{colim}}\nolimits}
\def\Spec{\mathop{\mathrm{Spec}}}
\def\Hom{\mathop{\mathrm{Hom}}\nolimits}
\def\Ext{\mathop{\mathrm{Ext}}\nolimits}
\def\SheafHom{\mathop{\mathcal{H}\!\mathit{om}}\nolimits}
\def\SheafExt{\mathop{\mathcal{E}\!\mathit{xt}}\nolimits}
\def\Sch{\mathit{Sch}}
\def\Mor{\mathop{\mathrm{Mor}}\nolimits}
\def\Ob{\mathop{\mathrm{Ob}}\nolimits}
\def\Sh{\mathop{\mathit{Sh}}\nolimits}
\def\NL{\mathop{N\!L}\nolimits}
\def\CH{\mathop{\mathrm{CH}}\nolimits}
\def\proetale{{pro\text{-}\acute{e}tale}}
\def\etale{{\acute{e}tale}}
\def\QCoh{\mathit{QCoh}}
\def\Ker{\mathop{\mathrm{Ker}}}
\def\Im{\mathop{\mathrm{Im}}}
\def\Coker{\mathop{\mathrm{Coker}}}
\def\Coim{\mathop{\mathrm{Coim}}}

% Boxtimes
%
\DeclareMathSymbol{\boxtimes}{\mathbin}{AMSa}{"02}

%
% Macros for moduli stacks/spaces
%
\def\QCohstack{\mathcal{QC}\!\mathit{oh}}
\def\Cohstack{\mathcal{C}\!\mathit{oh}}
\def\Spacesstack{\mathcal{S}\!\mathit{paces}}
\def\Quotfunctor{\mathrm{Quot}}
\def\Hilbfunctor{\mathrm{Hilb}}
\def\Curvesstack{\mathcal{C}\!\mathit{urves}}
\def\Polarizedstack{\mathcal{P}\!\mathit{olarized}}
\def\Complexesstack{\mathcal{C}\!\mathit{omplexes}}
% \Pic is the operator that assigns to X its picard group, usage \Pic(X)
% \Picardstack_{X/B} denotes the Picard stack of X over B
% \Picardfunctor_{X/B} denotes the Picard functor of X over B
\def\Pic{\mathop{\mathrm{Pic}}\nolimits}
\def\Picardstack{\mathcal{P}\!\mathit{ic}}
\def\Picardfunctor{\mathrm{Pic}}
\def\Deformationcategory{\mathcal{D}\!\mathit{ef}}
\usepackage{cleveref}
\crefname{originalsection}{original section}{original sections}
\Crefname{originalsection}{Original section}{Original sections}
\makeatletter
\newlabel{section-loc-qcoh-flat-base-change}{{0GQG}{source}{}{sourcepointer-1}{}}
\newlabel{section-loc-qcoh-flat-base-change@cref}{{[originalsection][0][]0GQG}{[1][0][]source}{}{}{}}
\makeatother
\crefname{originalproposition}{original proposition}{original propositions}
\Crefname{originalproposition}{Original proposition}{Original propositions}
\makeatletter
\newlabel{proposition-loc-qcoh-flat-base-change}{{0771}{source}{}{sourcepointer-2}{}}
\newlabel{proposition-loc-qcoh-flat-base-change@cref}{{[originalproposition][0][]0771}{[1][0][]source}{}{}{}}
\makeatother
\crefname{originallemma}{original lemma}{original lemmas}
\Crefname{originallemma}{Original lemma}{Original lemmas}
\makeatletter
\newlabel{lemma-flat-pullback-quasi-coherent}{{076X}{source}{}{sourcepointer-3}{}}
\newlabel{lemma-flat-pullback-quasi-coherent@cref}{{[originallemma][0][]076X}{[1][0][]source}{}{}{}}
\makeatother
\crefname{originalreference}{original reference}{original references}
\Crefname{originalreference}{Original reference}{Original references}
\makeatletter
\newlabel{stacks-sheaves-lemma-compare-quasi-coherent}{{0GQD}{source}{}{sourcepointer-4}{}}
\newlabel{stacks-sheaves-lemma-compare-quasi-coherent@cref}{{[originalreference][0][]0GQD}{[1][0][]source}{}{}{}}
\makeatother
\crefname{originallemma}{original lemma}{original lemmas}
\Crefname{originallemma}{Original lemma}{Original lemmas}
\makeatletter
\newlabel{lemma-quasi-coherent-check-exact}{{0GQF}{source}{}{sourcepointer-5}{}}
\newlabel{lemma-quasi-coherent-check-exact@cref}{{[originallemma][0][]0GQF}{[1][0][]source}{}{}{}}
\makeatother
\title{Exactness of a complex of quasi-coherent modules detected by parasitic cohomology}
\author{The Stacks Project Authors}
\date{stacks-project:stacks-cohomology.tex:local-snapshot-20261012; GNU FDL 1.2 or later}
\begin{document}
\maketitle
\noindent\textit{Editorial note.} All delivered mathematical text is the author's own, copied verbatim from the pinned Stacks Project chapter file stacks-cohomology.tex; no statement or proof step was rewritten, reordered or removed.\par\medskip
\noindent\textit{Editorial note.} The counted span cites the lemmas recorded in assembly.external\_references with canonical tags as navigation pointers.\par\medskip
\noindent\textit{Allowed standard prerequisites (editorial).}\begin{itemize}
\item Comparison of module categories on the small etale site. Cohomology of Algebraic Stacks, Proposition stacks-cohomology-proposition-loc-qcoh-flat-base-change. Identifies Ker/Im computed in the two topologies.
\item Flat pullback of quasi-coherent modules is exact. Cohomology of Algebraic Stacks, Lemma stacks-cohomology-lemma-flat-pullback-quasi-coherent; Sheaves on Stacks, Lemma stacks-sheaves-lemma-compare-quasi-coherent. Exactness of the qc restriction composite.
\item Checking a quasi-coherent module is zero. Cohomology of Algebraic Stacks, Lemma stacks-cohomology-lemma-quasi-coherent-check-exact. Concludes E = 0 from vanishing on flat points.
\end{itemize}\medskip
\ifdefined\hypertarget\hypertarget{sourcepointer-1}{}\fi
\noindent\textit{Source reference.} Original section 0GQG: sources/stacks-chow/tags.txt, source lines 20861--20861. Cross-chapter standard prerequisite cited by the author inside the counted span; it is used as a pointer and its own proof is not part of this item.\par\medskip
\ifdefined\hypertarget\hypertarget{sourcepointer-2}{}\fi
\noindent\textit{Source reference.} Original proposition 0771: sources/stacks-chow/tags.txt, source lines 9186--9186. Cross-chapter standard prerequisite cited by the author inside the counted span; it is used as a pointer and its own proof is not part of this item.\par\medskip
\ifdefined\hypertarget\hypertarget{sourcepointer-3}{}\fi
\noindent\textit{Source reference.} Original lemma 076X: sources/stacks-chow/tags.txt, source lines 9182--9182. Cross-chapter standard prerequisite cited by the author inside the counted span; it is used as a pointer and its own proof is not part of this item.\par\medskip
\ifdefined\hypertarget\hypertarget{sourcepointer-4}{}\fi
\noindent\textit{Source reference.} Original reference 0GQD: sources/stacks-chow/tags.txt, source lines 20858--20858. Cross-chapter standard prerequisite cited by the author inside the counted span; it is used as a pointer and its own proof is not part of this item.\par\medskip
\ifdefined\hypertarget\hypertarget{sourcepointer-5}{}\fi
\noindent\textit{Source reference.} Original lemma 0GQF: sources/stacks-chow/tags.txt, source lines 20860--20860. Cross-chapter standard prerequisite cited by the author inside the counted span; it is used as a pointer and its own proof is not part of this item.\par\medskip
\begin{lemma}
\label{lemma-exact-sequence-quasi-coherent-parasitic-cohomology}
Let $\mathcal{X}$ be an algebraic stack. Let
$\alpha : \mathcal{F} \to \mathcal{G}$ and
$\beta : \mathcal{G} \to \mathcal{H}$
be maps in $\QCoh(\mathcal{O}_\mathcal{X})$ with
$\beta \circ \alpha = 0$. The following are equivalent:
\begin{enumerate}
\item in the abelian category $\QCoh(\mathcal{O}_\mathcal{X})$
the complex $\mathcal{F} \to \mathcal{G} \to \mathcal{H}$
is exact at $\mathcal{G}$,
\item $\Ker(\beta)/\Im(\alpha)$ computed in either
$\textit{Mod}(\mathcal{X}_\etale, \mathcal{O}_\mathcal{X})$ or
$\textit{Mod}(\mathcal{X}_{fppf}, \mathcal{O}_\mathcal{X})$
is parasitic.
\end{enumerate}
\end{lemma}

\begin{proof}
We have $\QCoh(\mathcal{O}_\mathcal{X}) \subset
\textit{LQCoh}^{fbc}(\mathcal{O}_\mathcal{X})$, see
Section \ref{section-loc-qcoh-flat-base-change}.
Hence $\Ker(\beta)/\Im(\alpha)$ computed in
$\textit{Mod}(\mathcal{X}_\etale, \mathcal{O}_\mathcal{X})$ or
$\textit{Mod}(\mathcal{X}_{fppf}, \mathcal{O}_\mathcal{X})$ agree, see
Proposition \ref{proposition-loc-qcoh-flat-base-change}.
From now on we will use the \'etale topology on $\mathcal{X}$.

\medskip\noindent
Let $\mathcal{E}$ be the cohomology of
$\mathcal{F} \to \mathcal{G} \to \mathcal{H}$
computed in the abelian category $\QCoh(\mathcal{O}_\mathcal{X})$.
Let $x : U \to \mathcal{X}$ be a flat morphism where $U$ is a scheme.
As we are using the \'etale topology, the restriction functor
$\textit{Mod}(\mathcal{X}_\etale, \mathcal{O}_\mathcal{X})
\to \textit{Mod}(U_\etale, \mathcal{O}_U)$ is exact.
On the other hand, by Lemma \ref{lemma-flat-pullback-quasi-coherent} and
Sheaves on Stacks, Lemma \ref{stacks-sheaves-lemma-compare-quasi-coherent}
the restriction functor
$$
\QCoh(\mathcal{O}_\mathcal{X}) \xrightarrow{x^*}
\QCoh((\Sch/U)_\etale, \mathcal{O}) \xrightarrow{{-}|_{U_\etale}}
\QCoh(U_\etale, \mathcal{O}_U)
$$
is exact too. We conclude that
$\mathcal{E}|_{U_\etale} = (\Ker(\beta)/\Im(\alpha))|_{U_\etale}$.

\medskip\noindent
If (1) holds, then $\mathcal{E} = 0$ hence $\Ker(\beta)/\Im(\alpha)$
restricts to zero on $U_\etale$ for all $U$ flat over $\mathcal{X}$ and this
is the definition of a parasitic module. If (2) holds, then
$\Ker(\beta)/\Im(\alpha)$ restricts to zero on $U_\etale$ for all $U$
flat over $\mathcal{X}$ hence $\mathcal{E}$ restricts to zero on
$U_\etale$ for all $U$ flat over $\mathcal{X}$. This certainly implies
that the quasi-coherent module $\mathcal{E}$ is zero, for example apply
Lemma \ref{lemma-quasi-coherent-check-exact}
to the map $0 \to \mathcal{E}$.
\end{proof}

\end{document}
